\section{Statistical realization of multilevel prediction bounds}
\label{sec:v30-multilevel}
The occupation theorem concerns local distribution bounds, whereas a
statistical prediction problem also requires a common decoder and causal
state transitions. We now identify a class on which this distinction can
be evaluated exactly. The class consists of regenerative experiments on
finitely many committed routes. Its analysis keeps the entire acquired
law, including its atoms. In particular, no event requiring a favorable
report at every checkpoint is introduced.

\subsection{Regenerative experiments and the retained resource}
Fix integers $R,J\ge1$. Before observing any data, a route
$r\in\{1,\ldots,R\}$ is selected using a tape-independent random seed.
The route is then fixed. At stage $j$ a fresh latent variable and its
acquisition block are independent of all other stages, conditionally on
$r$. Write $H_{rj}$ for that block. After its compression, an independent
query $u$ is drawn from a finite set $\mathcal U_{rj}$ with probability
$p_{rj}(u)>0$. Only this query is executed. Its conditional mean, given
the complete history, is $q_{rj}(H_{rj},u)\in[0,1]$. Independence makes
the current block sufficient for this conditional mean; it does not
make that block available to the decoder.

At every queried boundary the predictor retains at most $M$ labels.
A transition reads the preceding label and the current acquisition
block, but the decoder reads only the retained label, the query, the
stage, the announced route and the independent seed. The stage and
initial route are read-only experiment descriptors. Previous outputs,
commands and reports cannot be read back. Temporary arithmetic is not
charged, as in the persistent-label model of this article. A route
cannot be changed after observing a block. Randomness correlated with
the tape is not an admitted seed.

For each route and stage define the one-block distortion
\begin{equation}\label{eq:v30-one-block}
 d_{rj}(M)=\min_{c_1,\ldots,c_M\in[0,1]^{\mathcal U_{rj}}}
 \E\min_{1\le m\le M}
       \sum_{u\in\mathcal U_{rj}}p_{rj}(u)
       |q_{rj}(H_{rj},u)-c_m(u)|^2.
\end{equation}
Repeated centers are allowed. The minimum exists: the domain is compact,
and the integrand is bounded and continuous in the centers. The
statistical objective, denoted $\mathcal R_M^{\rm reg}$, is the infimum
over the initial route law and one causal predictor of the maximum,
over $1\le j\le J$, of the expected query-averaged excess at stage $j$.
The expectation includes the initial seed. Thus the maximum is outside
the seed average, not inside it.

\begin{theorem}[Exact multilevel statistical minimax]
\label{thm:v30-regenerative}
Under the preceding protocol, for every integer $M\ge1$,
\begin{equation}\label{eq:v30-exact-lp}
 \mathcal R_M^{\rm reg}
 =\min_{\rho\in\Delta_R}\max_{1\le j\le J}
                  \sum_{r=1}^R\rho_r d_{rj}(M)
 =\max_{\lambda\in\Delta_J}\min_{1\le r\le R}
                  \sum_{j=1}^J\lambda_j d_{rj}(M),
\end{equation}
where $\Delta_n$ is the probability simplex. All extrema are attained.
The first equality is realized by a single $M$-label causal predictor,
not by an oracle granting the decoder the preceding history.
\end{theorem}
\begin{proof}
Condition on the complete independent seed. The initial route is now
fixed, and at stage $j$ the decoder has at most $M$ prediction vectors.
Its state may depend on earlier blocks, but every realized vector is
one of these same $M$ vectors. For each current block its loss is
therefore at least the distance to the nearest of those vectors.
The current block has its prescribed law after the seed is fixed.
Integrating proves a conditional lower bound $d_{rj}(M)$ at every
stage. This argument allows arbitrary randomized assignments, since
an average of distances cannot be smaller than their minimum.
If $\rho_r$ is the probability of the initial route, averaging at each
fixed stage gives a lower risk vector with coordinates
$\sum_r\rho_r d_{rj}(M)$. Taking the checkpoint maximum only now,
and then the infimum over predictors and route laws, proves the first
lower inequality in \eqref{eq:v30-exact-lp}.

Conversely, choose a minimizing $\rho$ and, for every $(r,j)$, a
minimizing codebook in \eqref{eq:v30-one-block}. On reading the current
block, the transition replaces the old label by the index of a nearest
center; ties are resolved by the least index. It retains no other
information. The decoder returns that center's coordinate at the
independently sampled query. The rule at a later stage again reads its
fresh block and replaces the label. These stage-indexed maps constitute
one causal predictor with a fixed program and the same label set at
all stages. They attain every conditional distortion $d_{rj}(M)$,
so their risk is the first minimum in \eqref{eq:v30-exact-lp}.
The finite-dimensional objective is continuous on $\Delta_R$.

For completeness, let $K$ be the convex hull of the row vectors
$(d_{r1}(M),\ldots,d_{rJ}(M))$ and put
$t_* = \min_{y\in K}\max_j y_j$. The point $t_*\mathbf1$ is on
the boundary of the closed convex upward set $K+\R_+^J$. A supporting
hyperplane has a nonzero normal $\lambda\ge0$; normalize it so that
$\sum_j\lambda_j=1$. Its support value is
$\min_{y\in K}\lambda\cdot y=t_*$, which equals the minimum over
the finitely many row vectors. Conversely, for every such $\lambda$
and every $y\in K$, $\lambda\cdot y\le\max_jy_j$.
These two inequalities prove the dual equality and its attainment.
This is finite convex separation, not a new linear-programming duality
principle; compare \cite{Rockafellar1964}.
\end{proof}

The theorem permits arbitrary finite one-block query menus and arbitrary
acquired laws. Its realizability assertion uses regeneration and initial
commitment essentially. It does not describe a later data-dependent
route change, a nonregenerative posterior, or a decoder with unspecified
side information. If the announced route must instead be retained in
the state, the construction above uses at most $RM$ labels, or at most
$\lceil\log_2R\rceil$ additional bits. No equality for that different
charged-route optimum is inferred from \eqref{eq:v30-exact-lp}.

\subsection{A positive experiment with an exact common design factor}
At every stage use the two-trial positive detector from
Section~\ref{sec:v27-evaluated}, with an independent uniform latent
variable on $[0,1]$. Its first command is uniform on $[1/4,3/4]^2$,
its cells are $k_0(t)=1/2+t/4$ and $k_1(t)=1/2-t/4$, and all first
reports are retained in the law. The acquired posterior mean $Z$ has
law
\begin{equation}\label{eq:v30-law}
 \nu=\frac5{16}\delta_{8/15}+\frac3{16}\delta_{4/9}+f(z)\,dz.
\end{equation}
The continuous component has mass $1/2$, is supported on
$[10/21,14/27]$, and has density bounded by $H=26$. These facts,
including the report evidence, are proved in
Proposition~\ref{prop:v27-density}. The four original query means are
\[
 q(z)=(3/4,\ 1/2+z/8,\ 1/2-z/8,\ 1/4),\qquad
 \|q(z)-q(z')\|_q^2=(z-z')^2/128.
\]
Here the four queries have equal weight.

Fix a matrix of gains $s=(s_{rj})$ with $0<s_{rj}\le1$. On route $r$
at stage $j$, execute the original second trial for its sampled query.
The scored binary event equals its failure indicator with independent
probability $s_{rj}$, and otherwise equals an independent fair coin.
This is a specified randomization of the readout, not a change to the
acquired law or an unobserved selection of histories. The conditional
query vector is
\begin{equation}\label{eq:v30-attenuation}
 q_{rj}(z)=(1-s_{rj})\mathbf1/2+s_{rj}q(z).
\end{equation}
It is an admitted probability vector and its squared query distance is
$s_{rj}^2(z-z')^2/128$. Every run executes exactly $2J$ detector trials;
the independent readout coins are not additional observations of a
latent variable. The entire protocol, including the restriction to a
committed route, is fixed before data are generated.

Put $a_{rj}=s_{rj}^2$ and define the finite design value
\begin{equation}\label{eq:v30-design}
 V(a)=\min_{\rho\in\Delta_R}\max_j\sum_r\rho_r a_{rj}
     =\max_{\lambda\in\Delta_J}\min_r\sum_j\lambda_j a_{rj}.
\end{equation}
It is positive. As before, $D_M(\nu)$ is scalar squared quantization
error for the full measure $\nu$, without deleting its atoms.

\begin{theorem}[Exact realization at every budget]
\label{thm:v30-positive-realization}
For the committed-route positive experiments just defined,
\begin{equation}\label{eq:v30-positive-risk}
                \mathcal R_M^{\rm reg}
                      =\frac{D_M(\nu)}{128}\,V(a)
                         \qquad(M\ge1).
\end{equation}
One initial route law attains this formula simultaneously for all
budgets, with a predictive codebook chosen anew at each budget.
Furthermore,
\begin{equation}\label{eq:v30-positive-asymptotic}
 \lim_{M\to\infty}M^2\mathcal R_M^{\rm reg}
 =\kappa_{\rm ex}V(a),\qquad
 \kappa_{\rm ex}=\frac1{1536}
               \left(\int_{10/21}^{14/27}f(z)^{1/3}\,dz\right)^3.
\end{equation}
\end{theorem}
\begin{proof}
Project arbitrary decoder vectors orthogonally onto the affine line
\eqref{eq:v30-attenuation} in the query norm. Projection does not
increase distance to a true vector on that line. Clipping the resulting
scalar centers to $[4/9,8/15]$, the convex hull of the support of $\nu$,
further decreases distance and yields admitted probability predictions.
The line isometry therefore gives
$d_{rj}(M)=s_{rj}^2D_M(\nu)/128$, exactly as in
Proposition~\ref{prop:v27-exact-reduction}. Substitution in
Theorem~\ref{thm:v30-regenerative} proves \eqref{eq:v30-positive-risk}.
A minimizing route law for $V(a)$ is independent of $M$. A common
scalar quantizer may be used on every route and stage, with the
appropriate readout gain, so the construction does not require a new
acquisition law at each budget. Finally,
Lemma~\ref{lem:v27-quantization} applies to the actual density and its
two atoms and proves \eqref{eq:v30-positive-asymptotic}.
\end{proof}

\subsection{Atom-aware capacities and a uniform relaxation gap}
Represent the committed routes by $R$ disjoint source--sink paths, each
with $J$ internal vertices $(r,j)$. Let $L_j$ be the query-averaged
squared excess, after clipping predictions to $[0,1]$. Thus $0\le L_j\le1$.
Use the whole probability space as the event in the occupation theorem,
so its event mass is one. Define
\begin{equation}\label{eq:v30-capacity}
 c_{rj}(t)=\min\{1,\tfrac12+b_{rj}\sqrt t\},\qquad
 b_{rj}=\frac{2HLM}{s_{rj}},\quad H=26,\quad L=8\sqrt2.
\end{equation}
The term $1/2$ is the total mass of the two acquired atoms, not a
conditional probability after choosing a route. It is indispensable.

\begin{lemma}[Local capacities from the full acquired law]
\label{lem:v30-atom-capacity}
The functions \eqref{eq:v30-capacity} satisfy the local occupation
inequalities both unconditionally and seed by seed, with common
constants. Their terminal capacities are one. For $0\le x\le1$,
\begin{equation}\label{eq:v30-atom-cost}
 \phi_{rj}(x)=\int_0^1[x-c_{rj}(t)]_+\,dt
            =\frac{(x-1/2)_+^3}{3b_{rj}^2}.
\end{equation}
In particular,
\begin{equation}\label{eq:v30-unit-cost}
                 \phi_{rj}(1)
                      =\frac{s_{rj}^2}{8306688\,M^2}.
\end{equation}
\end{lemma}
\begin{proof}
Fix a seed and a route. The decoder has at most $M$ query vectors.
The recovery map on the original line has norm $L$; after attenuation
its norm is $L/s_{rj}$. Therefore a query ball of radius $\sqrt t$
around an arbitrary vector restricts the acquired scalar to an interval
of length at most $2L\sqrt t/s_{rj}$. For $M$ vectors, the continuous
part has probability at most $2HLM\sqrt t/s_{rj}$. The two atoms have
total probability at most $1/2$, regardless of the centers. Adding
these bounds and clipping at one proves \eqref{eq:v30-capacity} on
the visited route. On every other route the visitation event is empty.
Averaging independent seeds preserves the same inequality. No
conditioning on the continuous component, on a report word at other
stages, or on an adaptively selected scalar is used.

Since $b_{rj}\ge 2HL>1$, $c_{rj}(1)=1$. For $x\le1/2$ the
integrand in \eqref{eq:v30-atom-cost} vanishes. For $x>1/2$, put
$y=x-1/2$. The integrand equals $y-b_{rj}\sqrt t$ from zero to
$(y/b_{rj})^2<1$, and zero thereafter. Its integral is
$y^3/(3b_{rj}^2)$. At $x=1$ this is
$1/(24b_{rj}^2)=s_{rj}^2/(96H^2L^2M^2)$.
The identity $96\cdot26^2\cdot128=8306688$ proves
\eqref{eq:v30-unit-cost}.
\end{proof}

\begin{theorem}[A statistically controlled occupation relaxation]
\label{thm:v30-controlled-gap}
Let $\Gamma$ and $\Gamma_{\rm pre}$ be respectively the occupation
and pre-acquisition perspective values of
Sections~\ref{sec:v29-occupation} and \ref{sec:v29-initial} for
\eqref{eq:v30-capacity}. Then
\begin{equation}\label{eq:v30-anchored-value}
 \Gamma_{\rm pre}
       =\frac{V(a)}{8306688\,M^2}
       \le\mathcal R_M^{\rm reg}
       \le\frac{8306688}{64800}\,\Gamma_{\rm pre}.
\end{equation}
The multiplicative upper bound is less than $129$ and is independent
of $R$, $J$, the positive gains and the budget. More precisely,
\begin{equation}\label{eq:v30-exact-gap}
 \frac{\mathcal R_M^{\rm reg}}{\Gamma_{\rm pre}}
            =\frac{8306688}{128}M^2D_M(\nu)
       \longrightarrow 8306688\kappa_{\rm ex}.
\end{equation}
The convergence is uniform over all finite route designs and positive
gain matrices. The scalar coefficient evaluation in
Section~\ref{sec:v29-evaluated} places the limiting ratio between
$4.05$ and $4.06$. If $R\ge2$, the unanchored value for these same
local capacities is $\Gamma=0$.
\end{theorem}
\begin{proof}
Conditioning on the initial independent seed fixes the first vertex
and hence the entire committed route. The anchored unit flow has
occupation one at every vertex of that route and zero elsewhere.
A mixture with route law $\rho$ therefore has level costs
$\sum_r\rho_r\phi_{rj}(1)$. Theorem~\ref{thm:v29-pre} and
\eqref{eq:v30-unit-cost} give
$\Gamma_{\rm pre}=V(a)/(8306688M^2)$ and its lower-bound direction.
Equivalently, the dual weights in \eqref{eq:v30-design} give the same
value; no nonoptimal primal value is used as a lower certificate.
Combining with \eqref{eq:v30-positive-risk} proves
\eqref{eq:v30-exact-gap}.

The support of $\nu$ is contained in an interval of length $4/45$.
Partition that entire interval into $M$ equal intervals and use their
midpoints. Every point of the support lies within $2/(45M)$ of a
center. Thus $D_M(\nu)\le4/(2025M^2)$, including the atoms.
Substitution in \eqref{eq:v30-exact-gap} gives the factor
$8306688/64800$ in \eqref{eq:v30-anchored-value}. The scalar
high-resolution limit proves the last equality in
\eqref{eq:v30-exact-gap}. Its right side has no route-design parameter,
so convergence is uniform even when the finite number of routes and
checkpoints, or the positive gains, vary with $M$.

For the last assertion, put mass $1/2$ on each of any two routes.
Every visited-node occupation is then at most $1/2$. Formula
\eqref{eq:v30-atom-cost} makes all node costs zero, so $\Gamma=0$.
In contrast $V(a)>0$, and hence $\Gamma_{\rm pre}>0$ for every finite
budget. This strict distinction arises in actual positive experiments,
not solely in freely assigned random-path loss models.
\end{proof}

\subsection{An unbounded loss from optimizing checkpoints separately}
For $J=R\ge2$ choose $s_{rj}=1$ if $r=j$ and $s_{rj}=1/J$ otherwise.
Every detector and acquisition law is the same; only the prescribed
readout gain along the committed route changes. The initial route
choice balances which checkpoint receives the larger gain.

\begin{corollary}[Exact multilevel consistency penalty]
\label{cor:v30-checkpoint-gap}
For this family,
\begin{equation}\label{eq:v30-balanced-family}
 V(a)=\frac1J+\frac1{J^2}-\frac1{J^3},\qquad
 \mathcal R_M^{\rm reg}
   =\frac{D_M(\nu)}{128}
           \left(\frac1J+\frac1{J^2}-\frac1{J^3}\right).
\end{equation}
The uniform route law and uniform checkpoint weights are matching
primal and dual optimizers, for every budget. Allowing a different
route law to minimize each checkpoint separately would instead give
$D_M(\nu)/(128J^2)$. The ratio of the actual optimum to that invalid
separate-checkpoint comparator is exactly $J+1-1/J$ for every finite
$M$.
\end{corollary}
\begin{proof}
At checkpoint $j$ the design cost under $\rho$ is
$J^{-2}+(1-J^{-2})\rho_j$. Its maximum is minimized when all
$\rho_j=1/J$, since $\max_j\rho_j\ge1/J$. This proves the first
formula. For uniform dual weights, every route has exactly the same
weighted cost, equal to the stated value, giving a matching dual.
At one separately optimized checkpoint one may choose a different
route and obtain gain squared $J^{-2}$. Division of the two exact
values proves the ratio. The denominator is positive because $\nu$
has a nonzero continuous part and no finite codebook gives it zero
distortion.
\end{proof}

The result is an unbounded penalty for dropping the common initial
route law, not an adaptivity advantage over admissible causal filters.
It applies to the explicitly attenuated regenerative query protocol,
not to the original tensor-only graph task. The two distinctions are
part of the theorem's information model. What the family supplies is
an exact statistical realization, a horizon-independent relaxation-gap
bound, and a multilevel dual witness for one common causal controller.
The scalar quantization and finite separation used in its proof are
classical; the earlier collision-uniform and nonregenerative graph
results remain separate, with their complete proofs.
